\documentclass[11pt]{article}
\usepackage[margin=1in]{geometry}
\usepackage{amsmath,amssymb,amsthm,mathtools}
\usepackage{graphicx,microtype,enumitem,xcolor}
\usepackage[colorlinks=true,linkcolor=blue!55!black,citecolor=blue!55!black,urlcolor=blue!55!black]{hyperref}

\newtheorem{theorem}{Theorem}
\newtheorem{lemma}{Lemma}
\newtheorem{remark}{Remark}
\newcommand{\B}{\mathbb B}
\newcommand{\Sph}{\mathbb S}
\newcommand{\R}{\mathbb R}
\newcommand{\Gammaf}{\Gamma}

\title{The Sharp Lipschitz Exponent for Ball-PSWF Heat Kernels\\When $\gamma\ge0$}
\author{Automated mathematics run \texttt{fa\_banach\_001}, lane 7}
\date{17 August 2026}

\begin{document}
\maketitle

\begin{center}
\fbox{\parbox{0.91\textwidth}{\textbf{Status.} Substantial partial result,
likely valid, pending human review.  The exponent in the source question is
proved to be the maximal value $\alpha=1$ for every $\gamma\ge0$.  The range
$-\tfrac12<\gamma<0$ remains open here.}}
\end{center}

\section{Source question and result}

Bonami, Kerkyacharian and Petrushev consider on the unit ball $\B^d$,
$d\ge2$, the measure
\[
 d\mu_\gamma(x)=(1-|x|^2)^{\gamma-1/2}\,dx,
 \qquad \gamma>-\tfrac12,
\]
and the positive ball-PSWF operator
\[
 L_c=-\Delta+\sum_{i,j=1}^d x_ix_j\partial_i\partial_j
          +(d+2\gamma)x\cdot\nabla+c^2|x|^2.
\]
Its intrinsic distance is
\[
 \rho(x,y)=\arccos\!\left(x\cdot y+sqrt{1-|x|^2}\sqrt{1-|y|^2}\right).
\]
Writing $p_t^c(x,y)$ for the heat kernel of $e^{-tL_c}$ and
$V(x,r)=\mu_\gamma(B_\rho(x,r))$, the source proves that some
$0<\alpha\le1$ satisfies its Gaussian H\"older estimate (8.23), and asks for
the dependence of $\alpha$ on $\gamma$ \cite[p.~40]{BKP2022}:

\begin{figure}[ht]
\centering
\includegraphics[width=\textwidth]{figures/open_problem_crop.png}
\caption{Equation (8.23) and the open question in \cite[p.~40]{BKP2022}.}
\end{figure}

\begin{theorem}[Nonnegative-parameter answer]\label{thm:main}
Let $d\ge2$, $\gamma\ge0$, and $c\ge0$.  There are constants $C,c_0>0$
such that, whenever $\rho(x,x')\le\sqrt t$,
\begin{equation}\label{eq:lip}
 |p_t^c(x,y)-p_t^c(x',y)|
 \le C\frac{\rho(x,x')}{\sqrt t}
 \frac{\exp\!\left(-\rho(x,y)^2/(c_0t)\right)}
 {[V(x,\sqrt t)V(y,\sqrt t)]^{1/2}}.
\end{equation}
Thus the exponent in the source question can be chosen to be
\[
                              \boxed{\alpha=1}
\]
throughout the range $\gamma\ge0$.
\end{theorem}

Since no H\"older exponent larger than $1$ is under consideration, this is
the sharp value of the exponent on this parameter range.

\section{Proof intuition}

The source's distance is not ad hoc: the map
\[
 x\longmapsto \bigl(x,s(x)\bigr),
 \qquad s(x)=\sqrt{1-|x|^2},
\]
identifies the ball with the upper hemisphere.  Under this map, the ball
measure becomes spherical measure weighted by $s^{2\gamma}$, and $-L_0$ is
the corresponding weighted Laplacian.

For $\gamma\ge0$ this weighted sphere has nonnegative Bakry--\'Emery
curvature.  The resulting reverse Poincar\'e inequality bounds the intrinsic
gradient of the heat kernel by $t^{-1/2}$ times a Gaussian kernel profile.
Integrating the gradient along the short spherical geodesic from $x$ to
$x'$ gives the Lipschitz factor $\rho(x,x')/\sqrt t$.  The source paper's
bounded-potential transfer theorem then carries exactly the same exponent
from $L_0$ to $L_c=L_0+c^2|x|^2$.

\section{The weighted hemisphere}

Let
\[
 H^d_+=\{(x,s)\in\Sph^d:s>0\},\qquad s=s(x)=\sqrt{1-|x|^2}.
\]
The pullback of spherical distance is precisely $\rho$.  Spherical volume in
the $x$ chart is $d\mathrm{vol}_{\Sph}=s^{-1}dx$, and hence
\begin{equation}\label{eq:measure}
 s^{2\gamma}\,d\mathrm{vol}_{\Sph}
       =(1-|x|^2)^{\gamma-1/2}\,dx=d\mu_\gamma(x).
\end{equation}
The weighted spherical diffusion is
\[
 \mathcal L_\gamma
   =\Delta_{\Sph^d}+2\gamma\,\nabla\log s\cdot\nabla.
\]
In the $x$ chart,
\[
 \mathcal L_\gamma
 =\sum_{i,j}(\delta_{ij}-x_ix_j)\partial_i\partial_j
      -(d+2\gamma)x\cdot\nabla=-L_0.
\]
Thus the unperturbed heat kernel from the source is the heat kernel of this
weighted hemisphere diffusion (with its natural weighted-Neumann closure).

Put $e^{-W}=s^{2\gamma}$, so $W=-2\gamma\log s$.  On the unit sphere,
\[
 \operatorname{Ric}=(d-1)g,
 \qquad \operatorname{Hess}(\log s)
       =-g-\frac{ds\otimes ds}{s^2}.
\]
Therefore its Bakry--\'Emery tensor is
\begin{equation}\label{eq:ricci}
 \operatorname{Ric}_W
 =\operatorname{Ric}+\operatorname{Hess}W
 =(d-1+2\gamma)g+2\gamma\frac{ds\otimes ds}{s^2}\ge0
\end{equation}
for every $\gamma\ge0$.  Equivalently, the diffusion satisfies
$\mathrm{CD}(0,\infty)$.  Formula \eqref{eq:ricci} can also be checked
directly on the polynomial core by the Bochner identity.  The associated
gradient inequality and its reverse-Poincar\'e consequence pass to the
closed Markov semigroup; see \cite[Chapters 3--4]{BGL2014}.  At $\gamma=0$
this is simply the ordinary spherical semigroup restricted to functions even
across the equator.

\section{From curvature and Gaussian bounds to Lipschitz bounds}

We isolate the short argument that retains the Gaussian localization.  Let
$P_t=e^{-tL_0}$, let $q_t$ be its heat kernel, and write
$\Gammaf(u)=|\nabla u|_g^2$.  The $\mathrm{CD}(0,\infty)$ reverse Poincar\'e
inequality is
\begin{equation}\label{eq:rp}
 2a\,\Gammaf(P_a f)
 \le P_a(f^2)-(P_af)^2\le P_a(f^2).
\end{equation}
Fix $t>0$ and $y$, put $a=t/2$, and apply \eqref{eq:rp} to
$f(z)=q_{t/2}(z,y)$.  Positivity and the semigroup identity give
\begin{align}
 \Gammaf_x(q_t(x,y))
 &\le \frac1t P_{t/2}\big(q_{t/2}(\cdot,y)^2\big)(x)\notag\\
 &\le \frac1t\|q_{t/2}(\cdot,y)\|_\infty q_t(x,y).
 \label{eq:kernel-rp}
\end{align}

The Gaussian upper estimate already proved in the source, together with
doubling, has the standard asymmetric consequences
\[
 \|q_{t/2}(\cdot,y)\|_\infty\le \frac{C}{V(y,\sqrt t)},
 \qquad
 q_t(x,y)\le \frac{C}{V(y,\sqrt t)}
       e^{-\rho(x,y)^2/(Ct)}.
\]
Taking square roots in \eqref{eq:kernel-rp} yields
\begin{equation}\label{eq:gradient}
 |\nabla_xq_t(x,y)|_g
 \le \frac{C}{\sqrt t\,V(y,\sqrt t)}
       e^{-\rho(x,y)^2/(Ct)}.
\end{equation}

Let $\delta=\rho(x,x')\le\sqrt t$ and integrate
\eqref{eq:gradient} along the minimizing spherical geodesic joining $x$ to
$x'$.  Along this segment, the triangle inequality preserves the Gaussian
factor after enlarging its denominator constant.  Hence
\[
 |q_t(x,y)-q_t(x',y)|
 \le C\frac{\delta}{\sqrt t\,V(y,\sqrt t)}
       e^{-\rho(x,y)^2/(Ct)}.
\]
Doubling gives, for some $Q>0$,
\[
 V(x,r)\le C\left(1+\frac{\rho(x,y)}r\right)^Q V(y,r).
\]
The polynomial factor is absorbed by the Gaussian, converting the last
display to
\begin{equation}\label{eq:q-lip}
 |q_t(x,y)-q_t(x',y)|
 \le C\frac{\rho(x,x')}{\sqrt t}
 \frac{e^{-\rho(x,y)^2/(Ct)}}
 {[V(x,\sqrt t)V(y,\sqrt t)]^{1/2}}.
\end{equation}

Finally, $L_c=L_0+c^2|x|^2$ is a bounded nonnegative potential
perturbation.  The source's H\"older-transfer theorem
\cite[Theorem~2.2]{BKP2022} preserves the exponent in
\eqref{eq:q-lip}; its application in Section 8.3 is precisely what produces
equation (8.23).  It therefore gives \eqref{eq:lip} with $\alpha=1$, proving
Theorem~\ref{thm:main}. \qed

\section{Verification, limitations, and novelty boundary}

The proof is analytic and has no computational dependency.  The generator
calculation agrees term by term with the source's $L_0$.  The key sign is in
\eqref{eq:ricci}: both additional terms are nonnegative exactly when
$\gamma\ge0$.  The kernel substitution in \eqref{eq:kernel-rp} then uses only
positivity, the semigroup law, the source's Gaussian bound, and doubling.

For $-\tfrac12<\gamma<0$, the normal curvature term
$2\gamma\,ds\otimes ds/s^2$ tends to $-\infty$ at the equator.  Thus this
proof supplies no uniform reverse Poincar\'e estimate.  The regular branch of
the one-dimensional Bessel boundary model suggests that $\alpha=1$ may still
be true, but a global scale-invariant gradient estimate was not obtained.
Eight focused attempts, including parameter-shift and spectral-kernel routes,
are recorded in the attempt log.

Bounded searches on 17 August 2026 covered the run indexes, exact source
metadata and question wording, weighted-ball gradient and Lipschitz phrases,
and the published source's citation list.  Its sole OpenAlex-listed citing
paper is unrelated to this exponent.  No explicit later determination was
found.  Novelty confidence is moderate because the weighted curvature method
is standard once the hemisphere identification is made.

Human review is recommended, focusing on the Markov closure at the equator
for $0<\gamma<1/2$, the pointwise use of reverse Poincar\'e, and the final
application of the source's perturbation theorem.

\begin{thebibliography}{9}
\bibitem{BKP2022}
A. Bonami, G. Kerkyacharian and P. Petrushev,
\emph{Gaussian bounds for the heat kernel associated to prolate spheroidal
wave functions with applications},
Constr. Approx. 56 (2022), 517--567; arXiv:2108.12917.

\bibitem{BGL2014}
D. Bakry, I. Gentil and M. Ledoux,
\emph{Analysis and Geometry of Markov Diffusion Operators},
Grundlehren der mathematischen Wissenschaften 348, Springer, 2014.
\end{thebibliography}

\end{document}
